\par\addvspace{12pt}\Needspace{9\baselineskip}

\section{等差数列与等比数列}

\label{sol:sec:analysis:2}

% ITEM: analysis-arithmetic-recover
\solutionitem{analysis-arithmetic-recover}{例1.6}
等差数列满足\(a_3+a_7=20,S_{10}=110\)。求通项与前\(n\)项和，再求使\(S_n>1000\)成立的最小正整数\(n\)。
\solutionanswer
设首项\(a\)、公差\(d\)。\(a_3+a_7=2a+8d=20\)，而\(S_{10}=10a+45d=110\)。解得\(a=2,d=2\)，故\(a_n=2n,S_n=n(n+1)\)。条件化为\(n(n+1)>1000\)。因\(31\cdot32=992\)，\(32\cdot33=1056\)，最小正整数为\(n=32\)。


% ITEM: analysis-geometric-recover
\solutionitem{analysis-geometric-recover}{例1.7}
等比数列满足\(a_2=6,a_5=162\)，求首项、公比及前\(n\)项和。另取\(b_n=(-1/2)^{n-1}\)，写出其部分和，比较原数列与部分和数列的变化。
\solutionanswer
由\(a_2=a_1q=6\)、\(a_5=a_1q^4=162\)，得\(q^3=27\)，故\(q=3,a_1=2\)。因此\(a_n=2\cdot3^{n-1}\)，\(S_n=3^n-1\)，严格递增且无界。对\(b_n=(-1/2)^{n-1}\)，
\[T_n=\sum_{k=1}^n b_k=\frac{1-(-1/2)^n}{1+1/2}=\frac23\left(1-(-1/2)^n\right).\]
\(b_n\)符号交替且幅度递减，故不单调；而\(T_{n+1}-T_n=b_{n+1}\)，其部分和增减交替。部分和在\(2/3\)两侧交替，并趋近该值。
\par
